\documentclass[UTF8,12pt, a4paper, oneside]{ctexart}
\usepackage{amsmath, amsthm, amssymb, appendix, bm, graphicx, hyperref, mathrsfs,float,tabularx}
% 这里是导言区
\author{WaltYoung}
\date{\today}
\linespread{1.5}
\begin{document}
\vspace*{4\baselineskip}
\centerline{\fontsize{36}{40}\selectfont\kaishu\makebox[4em][s]{温州大学}}
\begin{figure}[h]
    \centering
    \includegraphics[scale=0.8]{温大LOGO.png}
\end{figure}
\vspace*{2\baselineskip}
\centerline{\fontsize{24}{30}\selectfont\heiti{}大报告}
\vspace*{4\baselineskip}
\centerline
{
\fontsize{14}{14}\selectfont%
\renewcommand{\arraystretch}{1.4}
\begin{tabular}{c@{~~}c}
\makebox[4em][s]{题目}:&\underline{\makebox[18em][c]{Title}}\\
\makebox[4em][s]{姓名}:&\underline{\makebox[18em][c]{WaltYoung}}\\
\makebox[4em][s]{学号}:&\underline{\makebox[18em][c]{ID numben}}\\
\makebox[4em][s]{班级}:&\underline{\makebox[18em][c]{Class}}\\
\makebox[4em][s]{任课老师}:&\underline{\makebox[18em][c]{Teacher}}\\
\end{tabular}
}
% \maketitle
\clearpage
\tableofcontents
\clearpage
\section{摘要}

\paragraph{关键词}
\section{引言}

\section{概述}

\section{国内外现状}

\subsection{A}

\subsubsection{B}

% \begin{table}[H]
%     \centering
%     \caption{表名}
%     \begin{tabularx}{0.8\textwidth}{|X|X|X|X|}
%     \hline
%     A & B & C & D\\
%     \hline
%     1 & 2 & 3 & 4 \\
%     \hline
%     \end{tabularx}
% \end{table}

% \begin{figure}[H]
%     \centering
%     \includegraphics[scale=0.8]{}
%     \caption{图名}
% \end{figure}

% \textbf{加粗字体}
% \begin{itemize}
%     \item[·] a
%     \item[·] b
% \end{itemize}

% \begin{sloppypar}
% 针对特定区域设置自动换行，两端对齐
% \end{sloppypar}

% \cite{索引} % 引用文献项，索引为\bibitem{}括号内的文字

% $\{ E_1, E_2, E_3, \dots , E_n , E_{n+1}\}$ % 行内数学公式 {𝐸1, 𝐸2, 𝐸3, ⋯ , 𝐸𝑛}

% $$ \{ E_1, E_2, E_3, \dots , E_n , E_{n+1}\} $$ % 居中数学公式 {𝐸1, 𝐸2, 𝐸3, ⋯ , 𝐸𝑛}

\section{参考文献}
\noindent{[1]}

\noindent{[2]}

\noindent{[3]}

\noindent{[4]}

\noindent{[5]}

\noindent{[6]}

\noindent{[7]}

% 其他参考文献的语法
% \begin{thebibliography}{99}
%     \bibitem{a}作者. \emph{文献}[M]. 地点:出版社,年份.
%     \bibitem{b}作者. \emph{文献}[M]. 地点:出版社,年份.
%     \bibitem{c}作者. \emph{文献}[M]. 地点:出版社,年份.
%     \bibitem{d}作者. \emph{文献}[M]. 地点:出版社,年份.
%     \bibitem{e}作者. \emph{文献}[M]. 地点:出版社,年份.
%     \bibitem{f}作者. \emph{文献}[M]. 地点:出版社,年份.
%     \bibitem{g}作者. \emph{文献}[M]. 地点:出版社,年份.
%     \bibitem{h}作者. \emph{文献}[M]. 地点:出版社,年份.
%     \bibitem{i}作者. \emph{文献}[M]. 地点:出版社,年份.
%     \bibitem{j}作者. \emph{文献}[M]. 地点:出版社,年份.
%     \bibitem{k}作者. \emph{文献}[M]. 地点:出版社,年份.
%     \bibitem{l}作者. \emph{文献}[M]. 地点:出版社,年份.
% \end{thebibliography}
\end{document}